\chapter{Related work}
In this chapter the most promising works on the topic of argument mining will be exposed in order to give an idea about the latest methods employed in this field.\\
Here two main sections will be presented: one about the latest developments in argument mining for political debates which is one of the most similar applications to the the one of the CIGAIA dataset and one which focuses on state-of-the-art approaches that have been successfully applied in other fields but, to the best of the researcher’s knowledge, have not yet been explored for this application.
\section{Argument mining on political debates}
Political debates are one of the most seen application for argument mining as they have less noise than the one contained in other applications such as Twitter or Reddit controversies while still keeping a high quantity of argumentative phrases.\\
The main models who are used to solve these task are encoder-only transformers which are used
 to model the majority of tasks in argument mining (AC, ARI and ARC especially). 
BERT and the models derived from it (i.e. BERT, XLNET, RoBERTa, DistilBERT, and ALBERT)\cite{BERTbas} are the state-of-the-art as they output an embedding. Mamma  mia.